\renewcommand{\teacherTag}{QUESTION 4 \textbullet\ 35 MARKS}
\chapter{The Car Suspension: Laplace Transform and Damping}

\begin{scenario}
In a simple model of a car's suspension, the height of the road surface $x(t)$ drives the height of the car body $y(t)$ through
\[y''+b\,y'+10\,y=10\,x,\]
with units chosen so that the spring constant is 10. The constant $b$ is the setting of the shock absorber (damper): a larger $b$ means a harder damper, a smaller $b$ a softer one. A soft damper gives a bouncy ride and a hard one a slow, stiff ride. The test is the car driving onto a step in the road, $x(t)=u(t)$ (Figure~2 of the paper).
\end{scenario}

\begin{center}
\begin{tikzpicture}[>=Stealth,scale=0.9]
\fill[pattern=north east lines,pattern color=slate](-1,-0.3)rectangle(4,0);\draw(-1,0)--(4,0)--(4,1)--(9,1);
\fill[pattern=north east lines,pattern color=slate](4,0)rectangle(9,1);\fill[pattern=north east lines,pattern color=slate](4,-0.3)rectangle(9,0);
\draw[thick,navy,rounded corners=2pt](0,0.55)rectangle(2.8,1.2);\draw[thick,navy](0.6,1.2)--(0.9,1.7)--(2.0,1.7)--(2.3,1.2);
\draw[thick,navy](0.5,0.3)circle(0.28);\draw[thick,navy](2.3,0.3)circle(0.28);\draw[->,thick](3.0,0.9)--(3.6,0.9);
\draw[dashed](4,-0.4)--(4,2.1)node[above,font=\scriptsize]{$t=0$};
\draw[<->](9.3,0)--node[right,font=\scriptsize]{1}(9.3,1);
\node[font=\scriptsize,left]at(-1,0.15){road $x(t)$};\node[font=\scriptsize,left]at(-0.05,0.9){$y(t)$};
\end{tikzpicture}
\end{center}

\begin{keyidea}Laplace-transform the equation with zero initial conditions: $H(s)=\dfrac{10}{s^2+bs+10}$. For the step, $Y(s)=\dfrac{10}{s(s^2+bs+10)}$. The poles of $H$ decide everything: real $\Rightarrow$ no oscillation, complex $\Rightarrow$ bounce. Since $H(0)=1$, every stable setting settles at $y(\infty)=1$, the height of the step.\end{keyidea}

%------------------------------------------------------------------ 4(a)
\begin{qcard}{Q4(a)}{Hard damper, $b=7$}{12}
\textbf{Solve} for the response $y(t)$ of the car body with the damper set hard, $b=7$, and find out whether the body ever rises above the height of the step (that is, above 1).
\end{qcard}
\begin{solution}
\step{Step 1 --- Transform.} $s^2+7s+10=(s+2)(s+5)$, so
\[Y(s)=\frac{10}{s(s+2)(s+5)}=\frac As+\frac B{s+2}+\frac C{s+5}.\]
\step{Step 2 --- Cover-up.}
\[A=\frac{10}{(2)(5)}=1,\qquad B=\frac{10}{(-2)(3)}=-\frac53,\qquad C=\frac{10}{(-5)(-3)}=\frac23 .\]
(Check: $A+B+C=0=y(0)$ \checkmark.)
\step{Step 3 --- Invert.}
\[\boxed{y(t)=\Big(1-\tfrac53e^{-2t}+\tfrac23e^{-5t}\Big)u(t)}\]
\step{Step 4 --- Does it overshoot?} Differentiate:
\[y'(t)=\tfrac{10}3e^{-2t}-\tfrac{10}3e^{-5t}=\tfrac{10}3e^{-2t}\big(1-e^{-3t}\big)>0\quad\text{for all }t>0 .\]
$y$ is strictly increasing and tends to $1$, so it approaches 1 from below and \emph{never exceeds it}. (Over-damped: two real poles, $-2$ and $-5$.)
\ans{$y(t)=1-\frac53e^{-2t}+\frac23e^{-5t}$ ($t\ge0$). Since $y'(t)>0$ for all $t>0$ and $y\to1$, the body never rises above the step.}
\end{solution}

%------------------------------------------------------------------ 4(b)
\begin{qcard}{Q4(b)}{Soft damper, $b=2$}{13}
\textbf{Solve} for the response $y(t)$ of the car body with the damper set soft, $b=2$, and find out the frequency at which the body bounces.
\end{qcard}
\begin{solution}
\step{Step 1 --- Poles.} $s^2+2s+10=(s+1)^2+9$: complex poles $s=-1\pm j3$.

\step{Step 2 --- Partial fractions.}
\[Y(s)=\frac{10}{s\,[(s+1)^2+9]}=\frac As+\frac{Bs+C}{(s+1)^2+9},\qquad A=\frac{10}{10}=1 .\]
Then $10=(s^2+2s+10)+(Bs+C)s\Rightarrow$ $s^2$: $1+B=0\Rightarrow B=-1$; $\ s^1$: $2+C=0\Rightarrow C=-2$. So
\[Y(s)=\frac1s-\frac{s+2}{(s+1)^2+9}=\frac1s-\frac{s+1}{(s+1)^2+3^2}-\frac13\cdot\frac{3}{(s+1)^2+3^2}.\]
\step{Step 3 --- Invert} with $e^{-at}\cos\omega t\pair\frac{s+a}{(s+a)^2+\omega^2}$, $e^{-at}\sin\omega t\pair\frac{\omega}{(s+a)^2+\omega^2}$:
\[\boxed{y(t)=\Big[1-e^{-t}\Big(\cos3t+\tfrac13\sin3t\Big)\Big]u(t)}\]
(Check: $y(0)=1-1=0$ \checkmark. Differentiating, $y'(t)=e^{-t}\big(\cos3t+\tfrac13\sin3t\big)-e^{-t}\big(\cos3t-3\sin3t\big)=\tfrac{10}3e^{-t}\sin3t$, so $y'(0)=0$ \checkmark --- the body starts at rest.)

\step{Step 4 --- Bounce frequency.} The oscillation is $\cos3t$ inside a decaying envelope $e^{-t}$: the \emph{damped} natural frequency is the imaginary part of the poles,
\[\omega_d=3\text{ rad/s},\qquad f_d=\frac{3}{2\pi}=0.477\text{ Hz},\qquad T_d=\frac{2\pi}3=2.09\text{ (time units)}.\]
The first peak is where $y'=0$, i.e.\ $\sin3t=0$, $t=\pi/3$: $\;y(\pi/3)=1+e^{-\pi/3}=1.351$, a 35\% overshoot above the step.
\ans{$y(t)=1-e^{-t}(\cos3t+\frac13\sin3t)$; the body bounces at $\omega_d=3$ rad/s ($f\approx0.48$ Hz), first peak $1.35$ at $t=\pi/3$.}
\end{solution}

%------------------------------------------------------------------ 4(c)
\begin{qcard}{Q4(c)}{Critical damping}{10}
The car body oscillates when the poles are complex and does not when they are real. \textbf{Examine} the poles for the smallest value of $b$ at which the oscillation disappears, so that the car body reaches the height of the step as fast as possible without going above it.
\end{qcard}
\begin{solution}
\step{Step 1 --- Poles as a function of $b$.}
\[s^2+bs+10=0\quad\Longrightarrow\quad s=\frac{-b\pm\sqrt{b^2-40}}{2}.\]
Complex (oscillation) when $b^2<40$; real (no oscillation) when $b^2\ge40$.

\step{Step 2 --- Boundary.} The smallest $b$ with real poles is
\[b^2=40\ \Rightarrow\ \boxed{b=2\sqrt{10}\approx6.32},\qquad s_{1,2}=-\frac b2=-\sqrt{10}\approx-3.16\ \text{(a double pole)}.\]
This is \emph{critical damping} (damping ratio $\zeta=\frac{b}{2\sqrt{10}}=1$).

\step{Step 3 --- Response.} $Y(s)=\dfrac{10}{s(s+\sqrt{10})^2}=\dfrac1s-\dfrac1{s+\sqrt{10}}-\dfrac{\sqrt{10}}{(s+\sqrt{10})^2}$, so
\[y(t)=1-\big(1+\sqrt{10}\,t\big)e^{-\sqrt{10}\,t},\qquad y'(t)=10\,t\,e^{-\sqrt{10}t}>0 .\]
Monotonic, never above 1. For $b<2\sqrt{10}$ the poles become complex and the body overshoots; for $b>2\sqrt{10}$ one pole moves towards the origin (e.g.\ $-2$ for $b=7$), making the response slower. So $b=2\sqrt{10}$ gives the fastest rise without overshoot.

\begin{center}
\begin{tikzpicture}
\begin{axis}[width=12cm,height=6cm,axis lines=left,xlabel={$t$},ylabel={$y(t)$},xmin=0,xmax=6,ymin=0,ymax=1.45,ticklabel style={font=\scriptsize},label style={font=\scriptsize},legend style={font=\scriptsize,draw=none,fill=white,fill opacity=0.85,text opacity=1,at={(0.98,0.04)},anchor=south east},grid=major,grid style={gray!20},samples=200]
\addplot[very thick,coral,domain=0:6]{1-exp(-x)*(cos(deg(3*x))+sin(deg(3*x))/3)};\addlegendentry{soft, $b=2$ (bounces at 3 rad/s)}
\addplot[very thick,teal,domain=0:6]{1-(1+sqrt(10)*x)*exp(-sqrt(10)*x)};\addlegendentry{critical, $b=2\sqrt{10}\approx6.32$}
\addplot[very thick,navy,dashed,domain=0:6]{1-5/3*exp(-2*x)+2/3*exp(-5*x)};\addlegendentry{hard, $b=7$}
\addplot[thin,slate,dashed,domain=0:6]{1};
\addplot[only marks,mark=*,mark size=1.6pt,coral]coordinates{(1.0472,1.3509)};
\node[font=\scriptsize,coral,anchor=west]at(axis cs:1.12,1.37){$1+e^{-\pi/3}=1.35$};
\end{axis}
\end{tikzpicture}
\end{center}

\begin{center}
\begin{tikzpicture}[scale=0.62,>=Stealth]
\draw[->](-6,0)--(1.2,0)node[right,font=\scriptsize]{$\sigma$};\draw[->](0,-3.6)--(0,3.6)node[above,font=\scriptsize]{$j\omega$};
\node[coral,font=\large]at(-1,3){$\times$};\node[coral,font=\large]at(-1,-3){$\times$};\node[font=\tiny,coral,left]at(-1.2,3){$-1+3j$ ($b=2$)};
\node[teal,font=\large]at(-3.162,0){$\times$};\node[font=\tiny,teal,above]at(-3.162,0.25){$-\sqrt{10}$ (double)};
\node[navy,font=\large]at(-2,0){$\times$};\node[navy,font=\large]at(-5,0){$\times$};\node[font=\tiny,navy,below]at(-2,-0.25){$-2$};\node[font=\tiny,navy,below]at(-5,-0.25){$-5$};
\draw[dashed,gray](0,0)circle(3.162);
\draw[->,thick,slate](-1.15,2.8)arc(108:170:2.95);
\node[font=\tiny,slate,align=center]at(-4.2,2.4){increasing $b$:\\poles slide along\\$\abs s=\sqrt{10}$, meet,\\then split on the axis};
\end{tikzpicture}
\end{center}
\ans{Oscillation disappears when $b^2-40\ge0$; the smallest such value is $b=2\sqrt{10}\approx6.32$ (critical damping), with a double real pole at $s=-\sqrt{10}\approx-3.16$ and $y(t)=1-(1+\sqrt{10}t)e^{-\sqrt{10}t}$ --- the fastest response with no overshoot.}
\end{solution}
